\section{Theoretical Insights}
\label{sec:analysis}

In this section, we analyze how locality and sensor placement shape what can be recovered under sparse sensing. Our goal is to characterize when locality-aware reconstruction is well aligned with the underlying dynamics, and how missed spatial dependencies create non-identifiability under longitudinal sensing.
\subsection{Local Structure in Spatio-Temporal Dynamics}

We begin by characterizing the class of spatio-temporal processes that FieldFormer is designed to represent effectively. At a high level, a process is \emph{local} if the state at a given location and time is primarily determined by nearby conditions in space and recent history in time, rather than by distant regions of the domain or by the entire global state. This notion of locality is a structural property of many physical systems, not an assumption introduced by FieldFormer. A formal definition in terms of finite-order differential operators is provided in Appendix~\ref{sec:definitions} (Definition~\ref{def:local_operator}).

In practice, locality appears in two complementary forms:
\begin{enumerate}[leftmargin=*,topsep=0pt]
    \item processes in which influence spreads gradually through space over time, and
    \item processes in which influence propagates along paths with finite speed.
\end{enumerate}

\noindent\textbf{Diffusion-like processes.}
In diffusion-dominated systems, information spreads smoothly through repeated local interactions. Each location evolves based on nearby spatial context, and the effect of distant regions emerges only after many incremental updates. Examples include heat conduction or pollutant dispersion in the absence of strong advection. In such systems, accurate reconstruction depends on access to a sufficiently broad local neighborhood: missing nearby sensors can substantially degrade prediction quality, even when temporal observations are dense.

\noindent\textbf{Propagation-like processes.}
In propagation-dominated systems, changes travel along preferred directions and reach downstream locations after a delay, rather than diffusing uniformly in all directions. Examples include fluid transport, traffic congestion waves, and wind-driven pollutant plumes. Here, the state at a location may be influenced by a compact set of upstream points, making anisotropic local neighborhoods effective when their geometry is captured correctly. Learning which spatial and temporal neighbors matter can therefore be more important than simply increasing neighborhood size.

In classical PDE theory, diffusion-like processes are typically modeled by parabolic equations, while propagation-dominated processes are described by hyperbolic equations. Many real-world systems combine both behaviors at different spatial and temporal scales. These locality structures explain why neighborhood-based models can reconstruct coherent sensor-space models when the sensor network sufficiently covers the relevant local domains of dependence. FieldFormer aligns with this principle by constructing adaptive local neighborhoods and learning how information should be aggregated across space and time. We next formalize this alignment through an expressivity result, followed by a miss-coverage lower bound that captures the limits imposed by sparse sensing.

\subsection{Expressivity Under Available Local Context}

We first establish that FieldFormer has sufficient representational capacity to approximate locally governed spatio-temporal dynamics when the relevant local stencil information is available. This result should be interpreted as an architectural compatibility statement. It shows that local transformer aggregation can represent finite-stencil PDE dynamics, although this setting is different from extreme sparse observations we later consider during empirical evaluation.

\begin{theorem}[Expressivity for Finite-Stencil Local Dynamics]
\label{prop:universal_pde}
Let $u:\mathbb{R}^{d}\times\mathbb{R}\to\mathbb{R}^q$ solve a parabolic or hyperbolic PDE
\[
F\bigl(z,u(z),\nabla u(z),\ldots,D^r u(z)\bigr)=0,\qquad z=(x,t),
\]
with a finite-order, local operator. Assume a consistent \emph{explicit} finite-difference scheme
with compact stencil $\mathcal{S}\subset\{-s_{\max},\ldots,s_{\max}\}^d\times\{0,\ldots,k\}$ and a
CFL-admissible time step $\Delta t$. Let $U_h$ denote the numerical solution generated by this scheme, and write its local update as
\[
U_h^{n+1}(\mathbf{i})=\psi(v_{\mathbf{i},n}),\qquad
v_{\mathbf{i},n}
=
\bigl(U_h^{n-\ell_j}(\mathbf{i}+\mathbf{s}_j)\bigr)_{j=1}^{N},
\]
where $\{(\mathbf{s}_j,\ell_j)\}_{j=1}^{N}$ is an ordering of $\mathcal{S}$ and $N=|\mathcal{S}|$. For any compact $K\subset\mathbb{R}^{d+1}$, define the compact set of stencil-value tuples
\[
\mathcal{V}_K
:=
\left\{
v_{\mathbf{i},n}\in(\mathbb{R}^{q})^{N}
:
(x_{\mathbf{i}},t_{n+1})\in K
\right\}.
\]
Then for any $\varepsilon>0$ there exists a FieldFormer with neighborhood size $m\ge|\mathcal{S}|$, hidden width $d'$, and depth $L$ such that
\[
\sup_{v\in\mathcal{V}_K}
\left\|
\psi(v)-\mathrm{FF}_{\theta}(E(v))
\right\|_2
<\varepsilon,
\]
where $E(v)$ denotes the FieldFormer tokenization of the ordered stencil tuple $v$.
\end{theorem}

A detailed proof of the theorem is presented in Appendix~\ref{sec:proofs}.

\begin{remark}[Implication: Alignment with Local Dynamics]
Theorem~\ref{prop:universal_pde} shows that FieldFormer can represent the local update structure induced by finite-stencil discretizations of parabolic and hyperbolic PDEs. It controls approximation of the discrete update map $\psi$, not convergence to the continuous PDE solution $u$. A continuous-solution statement would require an additional numerical-analysis argument, for example
\[
\|u-\hat u\|
\le
\|u-U_h\|
+
\|U_h-\hat u\|,
\]
where $\|u-U_h\|$ is a discretization error controlled by assumptions such as consistency, stability, regularity, and grid refinement. The intended claim here is architectural compatibility with local finite-stencil dynamics when the relevant stencil information is available; it does not remove the ambiguity introduced when sparse sensing fails to observe important local dependencies.
\end{remark}

\subsection{Missed-Support Non-Identifiability under Longitudinal Sensing}

Even when the underlying dynamics are representable by the model class, reconstruction accuracy is limited by what the sensor network reveals. In sparse sensing regimes, errors arise not only from model mismatch or optimization, but also from \emph{missed spatial dependencies}: components of the local dynamics that are never observed by any sensor. This is the key distinction between expressivity and recoverability. A model may be capable of representing the correct update rule, while the observations may still be insufficient to determine which local dependencies are active.

We focus on longitudinal sensing regimes, where observations are dense in time but sparse in space. In this setting, sensors are fixed at a small subset of spatial locations and provide long temporal traces. Temporal density helps estimate local temporal evolution at the observed sites, but it does not compensate for missing spatial support. Whether a local dependency is observed or missed is therefore determined primarily by spatial sensor placement.

\noindent\textbf{Locality Assumption.}
We assume that the spatio-temporal process is governed by local interactions, as defined in Appendix~\ref{sec:definitions} (Assumption~\ref{ass:locality}). This ensures that the state update at any location depends only on a finite neighborhood in space and time.

\noindent\textbf{Sampling Assumption.}
We consider a longitudinal sensing regime in which sensors are fixed at a subset of spatial locations and provide measurements across consecutive time steps. As a result, whether a local dependency is observed or missed depends on spatial sensor placement, not on temporal sampling density. We model sensor placement as selecting spatial locations uniformly at random from the set of locations influencing local dynamics. This assumption reflects common deployments such as environmental monitoring networks and is formalized in Appendix~\ref{sec:definitions} (Assumption~\ref{ass:sampling}).

\noindent\textbf{Influence Assumption.}
We assume that missed spatial support can create a genuine ambiguity in the local update. To capture this, we associate an \emph{influence weight} with each neighborhood element, with weights summing to one and reflecting relative importance. In longitudinal settings, missing a spatial location removes all of its associated temporal information, so the total influence of a spatial location is the aggregate of its temporal components. The formal set-level assumption states that, for any missed spatial set, there can be an observationally indistinguishable stencil completion whose output differs by an amount proportional to the missed influence mass. The formal statement appears in Appendix~\ref{sec:definitions} (Assumption~\ref{ass:influence}).

Together, these assumptions formalize the central limitation of sparse sensing: even if the model class is expressive enough, unobserved spatial dependencies can leave multiple plausible local completions that agree with all measurements. Thus, failures of global reconstruction under extreme sparsity are not merely optimization failures, but consequences of the information geometry imposed by the sensor network.

\begin{theorem}[Missed-Support Non-Identifiability Under Longitudinal Sampling]
\label{prop:avg_long}
Sample $m_x$ spatial locations uniformly without replacement from $\mathcal{S}_x$, let $I\subseteq\mathcal{S}_x$ be the sensed set and $I^c$ the missed set.
Under ~\eqref{eq:set_inf} (in Appendix~\ref{sec:definitions}, Assumption~\ref{ass:influence}), let $\hat u$ be any estimator using only the observed stencil block $V_{I\times\mathcal{T}}$. For fixed $I$, define the two completion risks
\[
\mathrm{risk}_I(V)
:=
\mathbb{E}\!\left[
\|\hat u(V_{I\times\mathcal{T}})-\psi(V)\|_2
\mid I
\right],
\]
and
\[
\mathrm{risk}_I(T_{I^c}V)
:=
\mathbb{E}\!\left[
\|\hat u(V_{I\times\mathcal{T}})-\psi(T_{I^c}V)\|_2
\mid I
\right],
\]
where $T_{I^c}V$ agrees with $V$ on the observed coordinates and differs only on the missed spatial support. Then there exists a constant $c>0$ such that
\[
\inf_{\hat u}\;
\mathbb{E}_I
\left[
\max\{\mathrm{risk}_I(V),\mathrm{risk}_I(T_{I^c}V)\}
\right]
\;\ge\;
c\,\mathbb{E}_I\bigl[\,A(I^c)\,\bigr]
\;=\;
c\,\Bigl(1-\frac{m_x}{N_x}\Bigr).
\]
\end{theorem}

A detailed proof of the theorem is presented in Appendix~\ref{sec:proofs}.

\begin{remark}[Implication: Sensor Coverage and Recoverability]
The theorem is a non-identifiability statement, not a universal error law for every data distribution. It says that sparse sensing can leave two plausible stencil completions that are indistinguishable from the observed data; if their updates differ by the missed influence mass, no estimator can be accurate on both. Since longitudinal sensors provide dense time traces at fixed locations, observing a spatial site reveals all of its associated temporal information, while missing a spatial site removes all of it. The identity $\mathbb{E}_I[A(I^c)]=1-m_x/N_x$ captures how much influential spatial support is missed under random placement. In practice, recoverability also depends on sensor placement, noise, forcing, boundary conditions, and the inductive bias of the reconstruction model.
\end{remark}

\subsection{Connection to Sensor-Space Modeling}

The above results clarify why sensor-space modeling is a more identifiable objective under extreme sparsity. The expressivity result shows that FieldFormer is well matched to locally structured dynamics when the relevant local context is available. The missed-support non-identifiability result shows that sparse sensing can remove important spatial dependencies, leaving multiple globally plausible completions consistent with the same observations. Together, these results imply that global field recovery away from observational support is inherently prior-dependent: different models may produce different globally plausible reconstructions while remaining consistent with the same sparse sensor trajectories.

Sensor-space modeling avoids overinterpreting these globally underconstrained regions. Instead of claiming exact recovery of the full latent field, it evaluates whether a method can accurately reconstruct or forecast measurements over the persistent sensor topology. FieldFormer is designed for this setting: its adaptive local neighborhoods exploit the local domains of dependence that remain observed, while its mesh-free coordinate representation allows the same model to operate over evolving sensor networks and irregular observation layouts.
